% Shared closure-principle panel. Included by papers that state the full
% axiom basis, immediately after THREE_AXIOM_BASIS.tex. This states a
% global principle on the axiom class, not a fourth axiom; every surface
% counting the basis counts exactly three axioms. Reviewed reference:
% docs/AXIOM_REFERENCE.md.

\paragraph{The closure principle.}
The axioms describe the observer screen. The global closure principle states
that the simulating and the simulated description are one system. Every quantity that
has both a construction-side reading and a readback-side reading must take
the same value once a typed bridge proves that the two readings denote one
invariant. The equality is forced by self-identity. Constructing that bridge
and the return map remains part of the physics. Existence, uniqueness, and
stability are separate determinacy tests on the resulting equation. The name
given to such an equation carries no mathematical weight: after the typed
identification, unequal readings would describe two systems rather than the
single self-referential universe.

The local screen-grain proposal
\[
P=\varphi+\sqrt{\pi}/A_T(P),
\]
with \(A_T(P)\) the inverse coupling returned by the declared source map,
has one interval-certified fixed point on its physical domain; its
laboratory attachment awaits a same-scheme hadronic spectral transport.
The displayed map's certified fixed point sits at $P=1.6309720959$ and
returns $A_T=136.9948352$; the certified gauge-width variant's fixed
point sits at $P=1.6309682414$ and returns $137.0356601$. Through the
same outer equation the measured 2022 CODATA coupling
$\alpha^{-1}=137.035999177(21)$ reads back the grain value
$P=1.6309682094$. The comparisons stand at $3.0\times10^{-4}$ relative
for the displayed map and at $2.5\times10^{-6}$ for the gauge-width
variant, as diagnostics of the declared source maps.

The direct global proposal is
\[
N=\log M_0(\mathfrak U_N).
\]
The complete declared finite source class does not select a unique cosmic value, so
any successful completion needs an additional source law. A separate
common-load branch starts from
\[
N_0=\pi\exp\!\left(\frac{6\pi}{P\,\alpha_U(P)}\right).
\]
On the finite collar branch, the declared total reserve expectation \(P/4\)
and six-class equidistribution give presence probability \(P/24\) for each
declared class. If one class is physically selected as the blocked event, its
scalar-weighted presence receipt is discharged, and that normalized
collar-survival factor is proved to act on the global capacity, the
corresponding candidate is
\[
N_{\mathrm{pres}}=N_0\left(1-\frac{P}{24}\right),
\qquad
\ln\frac{N_{\mathrm{pres}}}{\pi}
=
\frac{6\pi}{P\,\alpha_U(P)}
+\ln\left(1-\frac{P}{24}\right).
\]
The exponential alternative \(N_{\mathrm{Pois}}=N_0e^{-P/24}\) requires a
separate mean-count or continuum carrier. Neither global attachment is
derived. The inherited screen/electroweak bridge premises, physical
common-load bridge, physical \(\mathbb Z_6\) seam action, one-class reserve
selection, scalar-weighted receipt, reserve-to-capacity map, and
horizon-record identification are open. On the source-forward numerical
branch,
\[
N_0=3.5321\times10^{122},
\]
so the finite-presence candidate is
\(3.2921\times10^{122}\) and the exponential candidate is
\(3.3001\times10^{122}\). The full weighted Planck
base-\(\Lambda\)CDM chain gives the comparison coordinate
\(3.3129\times10^{122}\). The respective residuals are
\(-0.63\) percent and \(-0.39\) percent. Both comparisons were exposed before
the branch choice and carry no predictive weight.
